% 进展报告封面：A4 绝对定位，依据学校 Word 封面的位置与字号排版。
% 采用 TikZ 仅影响封面；字体由 fonts/font-setup.tex 统一选择。
\newcommand{\MidtermLatin}{\rmfamily}
\newsavebox{\MidtermEnglishNormalBox}
\newsavebox{\MidtermEnglishModerateBox}
\newsavebox{\MidtermEnglishCompactBox}
\newdimen\MidtermEnglishAvailableHeight
\newdimen\MidtermEnglishTitleTopY
\newdimen\MidtermEnglishTitleBottomLimitY
\newcommand{\MidtermPrepareEnglishTitle}[2]{%
  \sbox{#1}{{\MidtermLatin
    \parbox[c]{180mm}{%
      \centering\setlength{\parskip}{0pt}%
      \linespread{#2}\fontsize{22pt}{22pt}\selectfont
      \MidtermTitleEn}}}%
}


\newcommand{\MidtermCoverNode}[3]{%
  \node[anchor=center,inner sep=0pt] at
  ([xshift=#1mm,yshift=-#2mm]current page.north west) {#3};%
}
\newcommand{\MidtermCoverField}[3]{%
  \node[anchor=west,inner sep=0pt] at
    ([xshift=45.5mm,yshift=-#1mm]current page.north west)
    {{\MidtermSong\fontsize{15pt}{17pt}\selectfont\makebox[55mm][s]{#2}}};%
  \node[anchor=west,inner sep=0pt] at
    ([xshift=100.5mm,yshift=-#1mm]current page.north west)
    {{\MidtermSong\fontsize{15pt}{17pt}\selectfont\underline{\makebox[64mm][c]{#3}}}};%
}
\begin{titlepage}
\thispagestyle{empty}
\MidtermPrepareEnglishTitle{\MidtermEnglishNormalBox}{1.5}%
\MidtermPrepareEnglishTitle{\MidtermEnglishModerateBox}{1.35}%
\MidtermPrepareEnglishTitle{\MidtermEnglishCompactBox}{1.2}%
% Page-anchored cover drawing must remain in the main document.
\ifcsname tikzexternaldisable\endcsname\tikzexternaldisable\fi
\begin{tikzpicture}[remember picture,overlay]
  \MidtermCoverNode{105}{53}{%
    {\MidtermSong\fontsize{18pt}{25pt}\selectfont 西北农林科技大学}}
  \MidtermCoverNode{105}{64.5}{%
    {\MidtermSong\fontsize{18pt}{25pt}\selectfont
      \underline{\makebox[24mm][c]{\MidtermGrade}}级\MidtermDegree 研究生学位论文}}
  \MidtermCoverNode{105}{76}{%
    {\MidtermSong\fontsize{18pt}{25pt}\selectfont 进展报告}}

  % Word source: 22 pt titles with 1.5 multiple line spacing.
  % Set the local font baseline to 22 pt so linespread=1.5 applies once.
  % Center the Chinese title at its Word-reference position.
  % Place the English title relative to the Chinese text box, so the
  % separation stays consistent when the Chinese title has multiple lines.
  \node[anchor=center,inner sep=0pt] (midterm-title-zh) at
    ([xshift=105mm,yshift=-114mm]current page.north west) {%
      {\MidtermHei
       \parbox[c]{177mm}{%
         \centering\setlength{\parskip}{0pt}%
         \linespread{1.5}\fontsize{22pt}{22pt}\selectfont
         \MidtermTitleZh}}};
  % Preserve the 9.5 mm space between title groups. The English title
  % may use tighter internal leading only when necessary to leave room
  % above the first metadata field. The safe lower boundary is at 169 mm.
  \pgfextracty{\MidtermEnglishTitleTopY}{\pgfpointanchor{midterm-title-zh}{south}}
  \advance\MidtermEnglishTitleTopY by -9.5mm
  \pgfextracty{\MidtermEnglishTitleBottomLimitY}{\pgfpointanchor{current page}{north west}}
  \advance\MidtermEnglishTitleBottomLimitY by -169mm
  \MidtermEnglishAvailableHeight=\dimexpr
    \MidtermEnglishTitleTopY-\MidtermEnglishTitleBottomLimitY\relax

  \def\MidtermSelectedEnglishBox{\MidtermEnglishNormalBox}%
  \def\MidtermEnglishStretch{1.5}%
  \ifdim\dimexpr\ht\MidtermEnglishNormalBox+\dp\MidtermEnglishNormalBox\relax>\MidtermEnglishAvailableHeight
    \def\MidtermSelectedEnglishBox{\MidtermEnglishModerateBox}%
    \def\MidtermEnglishStretch{1.35}%
    \ifdim\dimexpr\ht\MidtermEnglishModerateBox+\dp\MidtermEnglishModerateBox\relax>\MidtermEnglishAvailableHeight
      \def\MidtermSelectedEnglishBox{\MidtermEnglishCompactBox}%
      \def\MidtermEnglishStretch{1.2}%
      \ifdim\dimexpr\ht\MidtermEnglishCompactBox+\dp\MidtermEnglishCompactBox\relax>\MidtermEnglishAvailableHeight
        \PackageError{nwafu-midterm}{English title exceeds available cover height}%
          {Shorten the English title; the Chinese-English separation is fixed.}%
      \fi
    \fi
  \fi
  \typeout{Midterm English title line stretch: \MidtermEnglishStretch}%
  \node[anchor=north,inner sep=0pt] (midterm-title-en) at
    ([yshift=-9.5mm]midterm-title-zh.south) {\usebox{\MidtermSelectedEnglishBox}};

  \MidtermCoverField{180.2}{学\hfill 院（系、所）}{\MidtermCollege}
  \MidtermCoverField{190.9}{一级学科／专业领域}{\MidtermMajor}
  \MidtermCoverField{201.7}{研究生姓名及学号}{\MidtermStudent\quad\MidtermStudentId}
  \MidtermCoverField{212.5}{指导教师姓名}{\MidtermSupervisor}
  \MidtermCoverField{223.3}{中期考核小组组长签名}{\MidtermChair}
  \MidtermCoverField{234.1}{中期考核小组成员签名}{\MidtermMembersFirst}
  \node[anchor=west,inner sep=0pt] at
    ([xshift=100.5mm,yshift=-244.9mm]current page.north west)
    {{\MidtermSong\fontsize{15pt}{17pt}\selectfont
      \underline{\makebox[64mm][c]{\MidtermMembersSecond}}}};

  \MidtermCoverNode{105}{263.1}{%
    {\MidtermSong\fontsize{15pt}{20pt}\selectfont
      研究生通过中期考核日期：\quad\MidtermDate}}
\end{tikzpicture}
\ifcsname tikzexternalenable\endcsname\tikzexternalenable\fi
\null
\end{titlepage}
